\section{Kesimpulan}

Pelaksanaan magang industri di PT Inti Utama Solusindo selama periode Januari
hingga Juni 2026 memberikan pengalaman kerja yang sangat berharga dan relevan
dengan bidang studi Teknik Informatika. Berdasarkan seluruh kegiatan yang telah
dilakukan, beberapa kesimpulan dapat dirumuskan sebagai berikut:

\begin{enumerate}
      \item Penerapan teori secara langsung di industri menghasilkan pemahaman yang jauh
            lebih mendalam dibandingkan pembelajaran di kelas. Konsep-konsep seperti SDLC,
            MVC, state management, dan API \textit{integration} yang dipelajari secara
            teoretis di perkuliahan terasa jauh lebih konkret ketika diterapkan pada produk
            nyata yang digunakan oleh pengguna secara luas.
      \item Pengerjaan produk HealthyOne dan Canvasser memberikan pemahaman baru mengenai
            kompleksitas pengembangan sistem informasi di sektor kesehatan, di mana akurasi
            tampilan dan keandalan data bukan sekadar nilai tambah, melainkan kebutuhan
            utama yang tidak bisa dikompromikan.
      \item Pendekatan \textit{Server-Driven UI} (SDUI) yang diterapkan perusahaan
            merupakan pengalaman baru yang jarang ditemui di lingkungan akademik.
            Pendekatan ini memberikan perspektif berbeda mengenai bagaimana keputusan
            arsitektur dapat secara signifikan memengaruhi cara tim \textit{front-end}
            bekerja dan berkolaborasi dengan tim \textit{back-end}.
      \item Pergeseran peran dari \textit{Front-End Developer} menjadi \textit{Full-Stack}
            merupakan perkembangan yang signifikan selama periode magang ini. Perubahan
            tersebut memberikan kesempatan untuk mengembangkan kemampuan yang lebih luas,
            mulai dari pengembangan antarmuka hingga implementasi logika bisnis dan
            integrasi dengan layanan \textit{back-end}. Pengalaman tersebut juga
            meningkatkan kemampuan adaptasi, pemecahan masalah, serta pemahaman yang lebih
            menyeluruh terhadap proses pengembangan perangkat lunak secara end-to-end.
\end{enumerate}